\chapter{Amnesia}

\section*{Definition}
Amnesia is partial or complete loss of past memories (retrograde amnesia), inability to form new memories (anterograde amnesia), or both. It is different from ordinary forgetting and may occur with or without other cognitive problems. Causes range from medications and head injury to seizures, stroke, infection, and neurodegenerative disease; the pattern and timing help guide evaluation.

\section*{When to See a Doctor}

\begin{itemize}
 \item Any sudden memory loss, especially with headache, focal neurologic signs, impaired consciousness, seizure, or difficulty speaking
 \item Amnesia after head trauma with loss of consciousness or expanding time of post-traumatic amnesia
 \item Progressive memory decline over weeks to months interfering with daily function
 \item Amnesia with fever, seizure, confusion, new psychiatric symptoms, or severe headache
 \item Memory loss triggered by medication or substance use
\end{itemize}

\section*{How It Develops}
New declarative memories depend on networks that include the hippocampus, medial temporal lobes, thalamus, and connected brain regions. Injury or dysfunction in these networks can disrupt memory encoding or retrieval. The effects vary: some people retain attention, language, procedural learning, and personal identity, while others have broader cognitive or consciousness changes.

\section*{Causes}
\begin{itemize}
 \item Transient amnestic syndromes: transient global amnesia (TGA), usually sudden anterograde amnesia with repetitive questioning and preserved identity that resolves within 24 hours after dangerous causes are excluded; transient epileptic amnesia (TEA), caused by temporal-lobe seizures
 \item Traumatic brain injury: Concussion (post-traumatic amnesia, both retrograde and anterograde, duration correlates with injury severity), contusional injury to the temporal poles, diffuse axonal injury
 \item Stroke: Posterior cerebral artery infarction involving the hippocampus and parahippocampus; thalamic stroke (paramedian or tuberothalamic artery territory); anterior choroidal artery infarction
 \item Hypoxic-ischemic injury: Cardiac arrest, carbon monoxide poisoning, strangulation -- hippocampus is selectively vulnerable to hypoxia
 \item Neurodegenerative disease: Alzheimer's disease (early and prominent episodic memory loss), frontotemporal dementia (semantic variant presents with loss of word and object meaning), limbic-predominant age-related TDP-43 encephalopathy (LATE)
 \item Infectious: Herpes simplex encephalitis (temporal lobe inflammation), limbic encephalitis (paraneoplastic or autoimmune, anti-GAD, anti-LGI1, anti-NMDA receptor)
 \item Nutritional-toxic: thiamine deficiency (Wernicke encephalopathy and Korsakoff syndrome, which can cause confusion, severe memory impairment, and sometimes confabulation), chronic alcohol-related cognitive impairment, and carbon monoxide exposure
 \item Drug-induced: Benzodiazepines (especially triazolam, midazolam), anticholinergics (scopolamine, diphenhydramine), Z-drugs (zolpidem), medical cannabis
 \item Surgical: Bilateral medial temporal lobe resection (classic H.M. case), unilateral temporal lobectomy for epilepsy when contralateral hippocampus is compromised
\end{itemize}

\section*{Related Symptoms}
\begin{itemize}
 \item Confabulation
 \item Disorientation to time and place (spared orientation to person in most cases)
 \item Repetitive questioning
 \item Apathy
 \item Personality change
 \item Executive dysfunction
 \item Aphasia
 \item Apraxia
 \item Agnosia
 \item Visuospatial impairment
\end{itemize}
\textbf{Important:} New memory loss still needs medical assessment even when other cognitive abilities seem normal. Conditions such as transient epileptic amnesia, medication effects, and autoimmune encephalitis may require targeted testing; specialized imaging is reserved for selected cases when routine evaluation is inconclusive.

\section*{Medical Evaluation}
\begin{itemize}
 \item History including mode of onset (acute vs. insidious), temporal gradient of retrograde amnesia, presence of confabulation, risk factors (alcohol, head trauma, seizures, family history).
 \item Cognitive screening: Montreal Cognitive Assessment (MoCA) with emphasis on delayed recall, orientation, and clock drawing.
 \item Targeted laboratory studies, often including blood count, electrolytes, glucose, kidney and liver function, TSH, and vitamin B12; thiamine, infection, and autoimmune testing depends on the history and examination.
 \item Brain imaging, usually CT for acute trauma or suspected bleeding and MRI when stroke, tumor, inflammation, or a structural memory disorder is suspected.
 \item EEG when seizure or non-convulsive status is possible. FDG-PET or other specialized imaging is reserved for selected suspected neurodegenerative cases.
 \item Lumbar puncture with cerebrospinal-fluid testing only when infection, autoimmune encephalitis, malignancy, or another inflammatory process is suspected and it is safe to perform.
\end{itemize}

\section*{Treatment Options}
\begin{itemize}
 \item Treat reversible causes: prompt parenteral thiamine for suspected Wernicke encephalopathy without delaying urgently needed glucose, acyclovir for suspected herpes simplex encephalitis, and specialist-directed treatment for tumors or autoimmune encephalitis.
 \item Cholinesterase inhibitors (donepezil, rivastigmine) may provide symptomatic benefit in Alzheimer's disease; memantine may be considered in moderate-to-severe disease under clinician supervision.
 \item Review and, when appropriate, gradually reduce or substitute medicines that impair memory, such as anticholinergics and benzodiazepines; do not stop dependence-forming medicines abruptly without medical advice.
 \item Cognitive rehabilitation includes external memory aids, errorless learning techniques, and environmental modifications; psychological support for both patient and caregivers is essential for progressive etiologies.
\end{itemize}

\section*{Self-Care and Home Management}
\begin{itemize}
 \item Use memory aids: notebooks, digital voice recorders, smartphone reminders, calendars, and labeled storage to compensate for encoding deficits.
 \item Establish a consistent daily routine with fixed locations for keys, medications, and important objects.
 \item Encourage family members to use simple, direct statements rather than open-ended questions to reduce patient distress.
 \item Avoid alcohol and seek medical and nutritional support if Wernicke--Korsakoff syndrome is suspected; supplements should not replace urgent treatment for suspected thiamine deficiency.
 \item Supervise medication adherence and avoid sedative-hypnotics that worsen memory function.
\end{itemize}

\section*{References}
\begin{thebibliography}{99}
\bibitem{amnesia:ref1} Squire LR, Wixted JT. ``The cognitive neuroscience of human memory since H.M.'' \textit{Annu Rev Neurosci}. 2011;34:259-288.
\bibitem{amnesia:ref2} Bartsch T, Butler C, et al. ``Transient global amnesia: a review.'' \textit{J Neurol}. 2019;266(6):1513-1521.
\bibitem{amnesia:ref3} Budson AE, Solomon PR. \textit{Memory Loss, Alzheimer\'s Disease, and Dementia}, 3rd ed. Elsevier, 2021.
\bibitem{amnesia:ref4} Miller BL, Boeve BF, eds. \textit{The Behavioral Neurology of Dementia}, 2nd ed. Cambridge University Press, 2016.
\bibitem{amnesia:ref5} UpToDate. ``Amnesia: causes and approach to evaluation.'' Wolters Kluwer, 2023.
\bibitem{amnesia:ref6} American Academy of Neurology. Practice guidance and review resources for transient global amnesia and memory disorders.
\end{thebibliography}
